% GENERATED FILE. Edit canonical lessons, then run npm run generate:books.
% canonical-source-hash: fnv1a64:e963bc11c4083d48
% canonical-lessons: KA-C474-tyajya, KA-C474-marubalake, KA-C474-vingadisu, KA-C474-tyajisu, KA-C474-saurasakti

\chapter{Waste and Recycling}
\label{ch:ka-waste-and-recycling}
\hlchaptermodality{474}

\begin{chapteropening}
\textbf{By the end of this chapter:} \emph{I can say waste matter, recycling, to sort, to forsake and solar power.}

Complete the last lesson of chapter 474: I can say waste matter, recycling, to sort, to forsake and solar power.
\end{chapteropening}

\section[\texorpdfstring{\kn{ತ್ಯಾಜ್ಯ} (tyājya)}{ತ್ಯಾಜ್ಯ (tyājya)}]{\kn{ತ್ಯಾಜ್ಯ} (tyājya) — waste matter; refuse}

\label{lesson:KA-C474-tyajya}



\begin{quote}
Before the new one: say the Kannada for soot, then the Kannada for carbon.
\end{quote}

\subsection*{You'll want to know: \kn{ತ್ಯಾಜ್ಯ}}
\textbf{\kn{ತ್ಯಾಜ್ಯ}} — \emph{tyājya} — \textquotedblleft{}waste matter; refuse\textquotedblright{}.

\begin{quote}
\kn{ಕಾರ್ಖಾನೆಯ} \kn{ತ್ಯಾಜ್ಯ} \kn{ನದಿಗೆ} \kn{ಹರಿಯುತ್ತಿತ್ತು}.
\end{quote}

\begin{quote}
\emph{kārkhāneya tyājya nadige hariyuttittu.}
\end{quote}

\begin{quote}
Waste from the factory was flowing into the river.
\end{quote}

\kn{ಘನ} \kn{ತ್ಯಾಜ್ಯ}, solid waste; \kn{ಕಸ} is the everyday word.

\begin{cousinweb}
\textbf{\kn{ತ್ಯಾಜ್ಯ}} (\emph{tyājya}) is borrowed from Sanskrit \emph{tyājya}, \textquotedblleft{}to be left, to be abandoned\textquotedblright{} (Monier-Williams): what is thrown away.
\end{cousinweb}

\begin{grammarlens}[title={What you've built}]
One more word for getting around.
\end{grammarlens}

\subsection*{Your turn}
\begin{itemize}
\raggedright
  \item \emph{Say it:} \emph{tyājya}
  \item \emph{Say it:} \emph{tyājya}, once more
  \item \emph{Say it:} \emph{kārkhāneya tyājya nadige hariyuttittu.} — Waste from the factory was flowing into the river.
  \item \emph{Recall:} say the Kannada for soot, then the Kannada for carbon, then say \emph{tyājya} again
\end{itemize}

\begin{tcolorbox}[breakable,colback=teal!4,colframe=teal!35!black,title={Before you move on}]
What does \kn{ತ್ಯಾಜ್ಯ} mean? (\textquotedblleft{}Waste matter; refuse\textquotedblright{}.) Say it once more.
\end{tcolorbox}
\section[\texorpdfstring{\kn{ಮರುಬಳಕೆ} (marubaḷake)}{ಮರುಬಳಕೆ (marubaḷake)}]{\kn{ಮರುಬಳಕೆ} (marubaḷake) — recycling; reuse}

\label{lesson:KA-C474-marubalake}



\begin{quote}
Before the new one: say the Kannada for carbon, then the Kannada for waste matter.
\end{quote}

\subsection*{You'll want to know: \kn{ಮರುಬಳಕೆ}}
\textbf{\kn{ಮರುಬಳಕೆ}} — \emph{marubaḷake} — \textquotedblleft{}recycling; reuse\textquotedblright{}.

\begin{quote}
\kn{ಪ್ಲಾಸ್ಟಿಕ್} \kn{ಬಾಟಲಿಗಳ} \kn{ಮರುಬಳಕೆ} \kn{ಸಾಧ್ಯ}.
\end{quote}

\begin{quote}
\emph{plāsṭik bāṭaligaḷa marubaḷake sādhya.}
\end{quote}

\begin{quote}
Plastic bottles can be recycled.
\end{quote}

\kn{ಮರು}-, again, plus \kn{ಬಳಕೆ}, use, the noun of \kn{ಬಳಸು}.

\begin{cousinweb}
\textbf{\kn{ಮರುಬಳಕೆ}} (\emph{marubaḷake}) joins the native prefix \emph{maṟu}, \textquotedblleft{}again\textquotedblright{} (DEDR 4766), and \emph{baḷake}, \textquotedblleft{}use\textquotedblright{}, the noun of \kn{ಬಳಸು}.
\end{cousinweb}

\begin{grammarlens}[title={What you've built}]
One more word for getting around.
\end{grammarlens}

\subsection*{Your turn}
\begin{itemize}
\raggedright
  \item \emph{Say it:} \emph{marubaḷake}
  \item \emph{Say it:} \emph{marubaḷake}, once more
  \item \emph{Say it:} \emph{plāsṭik bāṭaligaḷa marubaḷake sādhya.} — Plastic bottles can be recycled.
  \item \emph{Recall:} say the Kannada for carbon, then the Kannada for waste matter, then say \emph{marubaḷake} again
\end{itemize}

\begin{tcolorbox}[breakable,colback=teal!4,colframe=teal!35!black,title={Before you move on}]
What does \kn{ಮರುಬಳಕೆ} mean? (\textquotedblleft{}Recycling; reuse\textquotedblright{}.) Say it once more.
\end{tcolorbox}
\section[\texorpdfstring{\kn{ವಿಂಗಡಿಸು} (viṅgaḍisu)}{ವಿಂಗಡಿಸು (viṅgaḍisu)}]{\kn{ವಿಂಗಡಿಸು} (viṅgaḍisu) — to sort; to separate}

\label{lesson:KA-C474-vingadisu}



\begin{quote}
Before the new one: say the Kannada for waste matter, then the Kannada for recycling.
\end{quote}

\subsection*{You'll want to know: \kn{ವಿಂಗಡಿಸು}}
\textbf{\kn{ವಿಂಗಡಿಸು}} — \emph{viṅgaḍisu} — \textquotedblleft{}to sort; to separate\textquotedblright{}.

\begin{quote}
\kn{ನಾವು} \kn{ಮನೆಯ} \kn{ಕಸವನ್ನು} \kn{ಒಣ} \kn{ಮತ್ತು} \kn{ಹಸಿ} \kn{ಕಸ} \kn{ಅಂತ} \kn{ವಿಂಗಡಿಸುತ್ತೇವೆ}.
\end{quote}

\begin{quote}
\emph{nāvu maneya kasavannu oṇa mattu hasi kasa anta viṅgaḍisuttēve.}
\end{quote}

\begin{quote}
We sort the household rubbish into dry and wet waste.
\end{quote}

Past (I …): \kn{ವಿಂಗಡಿಸಿದೆ}, \emph{viṅgaḍiside}.

What is sorted takes -\kn{ಅನ್ನು}.

\begin{cousinweb}
\textbf{\kn{ವಿಂಗಡಿಸು}} (\emph{viṅgaḍisu}) is a Kannada verb whose origin is not settled; DEDR shows no clear source for it.
\end{cousinweb}

\begin{grammarlens}[title={What you've built}]
One more word for getting around.
\end{grammarlens}

\subsection*{Your turn}
\begin{itemize}
\raggedright
  \item \emph{Say it:} \emph{viṅgaḍisu}
  \item \emph{Say it:} \emph{viṅgaḍisu}, once more
  \item \emph{Say it:} \emph{nāvu maneya kasavannu oṇa mattu hasi kasa anta viṅgaḍisuttēve.} — We sort the household rubbish into dry and wet waste.
  \item \emph{Recall:} say the Kannada for waste matter, then the Kannada for recycling, then say \emph{viṅgaḍisu} again
\end{itemize}

\begin{tcolorbox}[breakable,colback=teal!4,colframe=teal!35!black,title={Before you move on}]
What does \kn{ವಿಂಗಡಿಸು} mean? (\textquotedblleft{}To sort; to separate\textquotedblright{}.) Say it once more.
\end{tcolorbox}
\section[\texorpdfstring{\kn{ತ್ಯಜಿಸು} (tyajisu)}{ತ್ಯಜಿಸು (tyajisu)}]{\kn{ತ್ಯಜಿಸು} (tyajisu) — to forsake; to give up}

\label{lesson:KA-C474-tyajisu}



\begin{quote}
Before the new one: say the Kannada for recycling, then the Kannada for to sort.
\end{quote}

\subsection*{You'll want to know: \kn{ತ್ಯಜಿಸು}}
\textbf{\kn{ತ್ಯಜಿಸು}} — \emph{tyajisu} — \textquotedblleft{}to forsake; to give up\textquotedblright{}.

\begin{quote}
\kn{ಅವನು} \kn{ಪ್ಲಾಸ್ಟಿಕ್} \kn{ಚೀಲಗಳನ್ನು} \kn{ತ್ಯಜಿಸಿದನು}.
\end{quote}

\begin{quote}
\emph{avanu plāsṭik cīlagaḷannu tyajisidanu.}
\end{quote}

\begin{quote}
He gave up plastic bags.
\end{quote}

Past (I …): \kn{ತ್ಯಜಿಸಿದೆ}, \emph{tyajiside}.

A formal word; what is given up takes -\kn{ಅನ್ನು}.

\begin{cousinweb}
\textbf{\kn{ತ್ಯಜಿಸು}} (\emph{tyajisu}) is built from the Sanskrit root \emph{tyaj}, \textquotedblleft{}to leave, to let go\textquotedblright{} (CDIAL 5984, \emph{tyajati}), with the verb ending -\emph{isu}.
\end{cousinweb}

\begin{grammarlens}[title={What you've built}]
One more word for getting around.
\end{grammarlens}

\subsection*{Your turn}
\begin{itemize}
\raggedright
  \item \emph{Say it:} \emph{tyajisu}
  \item \emph{Say it:} \emph{tyajisu}, once more
  \item \emph{Say it:} \emph{avanu plāsṭik cīlagaḷannu tyajisidanu.} — He gave up plastic bags.
  \item \emph{Recall:} say the Kannada for recycling, then the Kannada for to sort, then say \emph{tyajisu} again
\end{itemize}

\begin{tcolorbox}[breakable,colback=teal!4,colframe=teal!35!black,title={Before you move on}]
What does \kn{ತ್ಯಜಿಸು} mean? (\textquotedblleft{}To forsake; to give up\textquotedblright{}.) Say it once more.
\end{tcolorbox}
\section[\texorpdfstring{\kn{ಸೌರಶಕ್ತಿ} (sauraśakti)}{ಸೌರಶಕ್ತಿ (sauraśakti)}]{\kn{ಸೌರಶಕ್ತಿ} (sauraśakti) — solar power}

\label{lesson:KA-C474-saurasakti}



\begin{quote}
Before the new one: say the Kannada for to sort, then the Kannada for to forsake.
\end{quote}

\subsection*{You'll want to know: \kn{ಸೌರಶಕ್ತಿ}}
\textbf{\kn{ಸೌರಶಕ್ತಿ}} — \emph{sauraśakti} — \textquotedblleft{}solar power\textquotedblright{}.

\begin{quote}
\kn{ನಮ್ಮ} \kn{ಶಾಲೆ} \kn{ಸೌರಶಕ್ತಿಯಿಂದ} \kn{ನಡೆಯುತ್ತದೆ}.
\end{quote}

\begin{quote}
\emph{namma śāle sauraśaktiyinda naḍeyuttade.}
\end{quote}

\begin{quote}
Our school runs on solar power.
\end{quote}

\kn{ಸೌರ}, of the sun, plus \kn{ಶಕ್ತಿ}, power.

\begin{cousinweb}
\textbf{\kn{ಸೌರಶಕ್ತಿ}} (\emph{sauraśakti}) is a modern compound built from Sanskrit words: \emph{saura}, \textquotedblleft{}relating to the sun\textquotedblright{}, and \emph{śakti}, \textquotedblleft{}power, energy\textquotedblright{} (Monier-Williams).
\end{cousinweb}

\begin{grammarlens}[title={What you've built}]
One more word for getting around.
\end{grammarlens}

\subsection*{Your turn}
\begin{itemize}
\raggedright
  \item \emph{Say it:} \emph{sauraśakti}
  \item \emph{Say it:} \emph{sauraśakti}, once more
  \item \emph{Say it:} \emph{namma śāle sauraśaktiyinda naḍeyuttade.} — Our school runs on solar power.
  \item \emph{Recall:} say the Kannada for to sort, then the Kannada for to forsake, then say \emph{sauraśakti} again
\end{itemize}

\begin{tcolorbox}[breakable,colback=teal!4,colframe=teal!35!black,title={Before you move on}]
What does \kn{ಸೌರಶಕ್ತಿ} mean? (\textquotedblleft{}Solar power\textquotedblright{}.) Say it once more.
\end{tcolorbox}
